% =====================================================================
%  Страница «О конспекте»: лектор, источники, как читать.
%  Имя лектора задаётся макросом \lecturer в main.tex.
% =====================================================================
\section*{О конспекте}
\addcontentsline{toc}{section}{О конспекте}

Конспект лекций по математической статистике (лектор~\lecturer, осенний семестр
2026~года).

Конспект опирается на лекции и учебник~\cite{bure}, по которому построена
значительная часть курса. Места, где изложение взято из учебника или
дополнено по нему, отмечены ссылкой вида~\cite[гл.~13, §4]{bure}.
Обозначения в основном следуют лекциям и учебнику; их сводка~--- на следующей странице.

Материал, которого не было на лекции (доказательства и пояснения из учебника
или от составителя (или от Claude)), вынесен в зелёные блоки <<Дополнение>>. При подготовке
к экзамену их можно пропускать.

В конце некоторых лекций есть блок <<Что нужно знать к экзамену>>: список
определений и теорем и места, где легко ошибиться.

\paragraph{Благодарности.} Спасибо одногруппникам, которые делились своими
конспектами и помогали искать материалы, в частности, Ивану~Шитову и Максиму~Маслову.

\paragraph{Нашли ошибку?} Конспект неофициальный: его пишут студенты по ходу
семестра, лектор его не проверял, при подготовке использовались ИИ-инструменты.
Поэтому в нём возможны неточности и опечатки, и он не заменяет лекции и учебник:
при расхождении верными считайте их. Материал предоставляется <<как есть>>,
составители не несут ответственности за последствия его использования.
Если заметили ошибку или непонятное место~--- напишите составителю или
откройте issue в репозитории \url{https://github.com/atelitsin/matstat-notes}.
Свежая версия PDF всегда лежит там же.

\clearpage
